% Options for packages loaded elsewhere
\PassOptionsToPackage{unicode}{hyperref}
\PassOptionsToPackage{hyphens}{url}
\PassOptionsToPackage{dvipsnames,svgnames,x11names}{xcolor}
\documentclass[
  UTF8,11pt,a4paper]{ctexart}
\usepackage{xcolor}
\usepackage[margin=2.2cm]{geometry}
\usepackage{amsmath,amssymb}
\setcounter{secnumdepth}{5}
\usepackage{iftex}
\ifPDFTeX
  \usepackage[T1]{fontenc}
  \usepackage[utf8]{inputenc}
  \usepackage{textcomp} % provide euro and other symbols
\else % if luatex or xetex
  \usepackage{unicode-math} % this also loads fontspec
  \defaultfontfeatures{Scale=MatchLowercase}
  \defaultfontfeatures[\rmfamily]{Ligatures=TeX,Scale=1}
\fi
\usepackage{lmodern}
\ifPDFTeX\else
  % xetex/luatex font selection
\fi
% Use upquote if available, for straight quotes in verbatim environments
\IfFileExists{upquote.sty}{\usepackage{upquote}}{}
\IfFileExists{microtype.sty}{% use microtype if available
  \usepackage[]{microtype}
  \UseMicrotypeSet[protrusion]{basicmath} % disable protrusion for tt fonts
}{}
\makeatletter
\@ifundefined{KOMAClassName}{% if non-KOMA class
  \IfFileExists{parskip.sty}{%
    \usepackage{parskip}
  }{% else
    \setlength{\parindent}{0pt}
    \setlength{\parskip}{6pt plus 2pt minus 1pt}}
}{% if KOMA class
  \KOMAoptions{parskip=half}}
\makeatother
\usepackage{listings}
\setmonofont{Microsoft YaHei UI}
\setCJKmonofont{Microsoft YaHei UI}
\lstset{
  basicstyle=\small\ttfamily,
  columns=fullflexible,
  keepspaces=true,
  breaklines=true,
  breakatwhitespace=false,
  showstringspaces=false
}
\newcommand{\passthrough}[1]{#1}
\lstset{defaultdialect=[5.3]Lua}
\lstset{defaultdialect=[x86masm]Assembler}
\usepackage{longtable,booktabs,array}
\newcounter{none} % for unnumbered tables
\usepackage{calc} % for calculating minipage widths
% Correct order of tables after \paragraph or \subparagraph
\usepackage{etoolbox}
\makeatletter
\patchcmd\longtable{\par}{\if@noskipsec\mbox{}\fi\par}{}{}
\makeatother
% Allow footnotes in longtable head/foot
\IfFileExists{footnotehyper.sty}{\usepackage{footnotehyper}}{\usepackage{footnote}}
\makesavenoteenv{longtable}
\setlength{\emergencystretch}{3em} % prevent overfull lines
\providecommand{\tightlist}{%
  \setlength{\itemsep}{0pt}\setlength{\parskip}{0pt}}
\usepackage{bookmark}
\IfFileExists{xurl.sty}{\usepackage{xurl}}{} % add URL line breaks if available
\urlstyle{same}
\hypersetup{
  colorlinks=true,
  linkcolor={teal},
  filecolor={Maroon},
  citecolor={Blue},
  urlcolor={blue},
  pdfcreator={LaTeX via pandoc}}

\title{Embedded Knowledge Base}
\usepackage{etoolbox}
\makeatletter
\providecommand{\subtitle}[1]{% add subtitle to \maketitle
  \apptocmd{\@title}{\par {\large #1 \par}}{}{}
}
\makeatother
\subtitle{ESP32 Draggy Racing Logger 工程知识库}
\author{Workspace Engineering Notes}
\date{2026-06-10}

\begin{document}
\maketitle

{
\setcounter{tocdepth}{3}
\tableofcontents
}
\section{Embedded Knowledge Base}\label{embedded-knowledge-base}

\begin{quote}
面向未来自己的工程复盘。本文以当前工作区源码为事实基线；旧
README、优化报告和面试笔记只作为设计背景，不覆盖代码事实。

最后核查：2026-06-10
\end{quote}

\subsection{目录}\label{ux76eeux5f55}

\begin{enumerate}
\def\labelenumi{\arabic{enumi}.}
\tightlist
\item
  \hyperref[ux7b2cux4e00ux90e8ux5206ux9879ux76eeux5168ux666fux56fe]{第一部分：项目全景图}
\item
  \hyperref[ux7b2cux4e8cux90e8ux5206ux7cfbux7edfux67b6ux6784ux8bbeux8ba1]{第二部分：系统架构设计}
\item
  \hyperref[ux7b2cux4e09ux90e8ux5206ux9879ux76eeux8fd0ux884cux6d41ux7a0b]{第三部分：项目运行流程}
\item
  \hyperref[ux7b2cux56dbux90e8ux5206ux6838ux5fc3ux6a21ux5757ux89e3ux6790]{第四部分：核心模块解析}
\item
  \hyperref[ux7b2cux4e94ux90e8ux5206ux6280ux672fux77e5ux8bc6ux63d0ux53d6]{第五部分：技术知识提取}
\item
  \hyperref[ux7b2cux516dux90e8ux5206ux8bbeux8ba1ux51b3ux7b56ux5206ux6790]{第六部分：设计决策分析}
\item
  \hyperref[ux7b2cux4e03ux90e8ux5206ux8c03ux8bd5ux4e0eux8e29ux5751ux8bb0ux5f55]{第七部分：调试与踩坑记录}
\item
  \hyperref[ux7b2cux516bux90e8ux5206ux5de5ux7a0bux5b9eux8df5ux603bux7ed3]{第八部分：工程实践总结}
\item
  \hyperref[ux7b2cux4e5dux90e8ux5206ux9879ux76eeux590dux76d8]{第九部分：项目复盘}
\item
  \hyperref[ux7b2cux5341ux90e8ux5206ux4e2aux4ebaux77e5ux8bc6ux5730ux56fe]{第十部分：个人知识地图}
\end{enumerate}

\begin{center}\rule{0.5\linewidth}{0.5pt}\end{center}

\subsection{阅读前先知道的事实}\label{ux9605ux8bfbux524dux5148ux77e5ux9053ux7684ux4e8bux5b9e}

这个仓库只有一个主固件项目，但它经历了明显演化：

\begin{enumerate}
\def\labelenumi{\arabic{enumi}.}
\tightlist
\item
  最初是 ESP32 + MPU6050 + OLED 的双 I2C 演示。
\item
  后来加入 GT-U7 GPS，目标变成低成本 Dragy 风格车辆性能仪。
\item
  再加入摩托车倾角/G 力分析、Wi-Fi 网页、音频录制、SD 卡日志、PCB 和 3D
  打印外壳。
\end{enumerate}

因此旧文档与当前代码有冲突。以后判断``现在是什么''，优先级应为：

\begin{lstlisting}
当前源码与 platformio.ini
        >
当前硬件接线和实测记录
        >
README / GPS 优化报告 / 面试笔记
\end{lstlisting}

当前已发现的版本漂移：

{\def\LTcaptype{none} % do not increment counter
\begin{longtable}[]{@{}lll@{}}
\toprule\noalign{}
主题 & 当前源码 & 旧文档中的其他说法 \\
\midrule\noalign{}
\endhead
\bottomrule\noalign{}
\endlastfoot
I2C & MPU: GPIO21/22；OLED: GPIO18/19，两条总线 & README 曾写 21/19 与
25/26；GPS\_WIRING 曾写共享 21/22 \\
GPS 串口 & UART2，GPIO16/17，9600 bps & 优化报告写过 38400 bps \\
GPS 更新率 & 启动配置 10 Hz；掉线恢复时改为 5 Hz & 文档同时出现 5 Hz、10
Hz \\
SD 卡 & 模块完整，但 \passthrough{\lstinline!setup()!} 初始化被注释 &
部分描述把它写成可直接使用 \\
GPS-IMU 融合 & \passthrough{\lstinline!loop()!} 每 5 ms 调用
\passthrough{\lstinline!fuseIMUSpeed()!} &
\passthrough{\lstinline!/data!} 的 \passthrough{\lstinline!imuSpeed!}
仍固定输出 0，注释称暂未实现 \\
摩托互补滤波 & 已在 \passthrough{\lstinline!performance\_analyzer.cpp!}
实现 & 面试笔记部分段落称姿态互补滤波未实现 \\
\end{longtable}
}

编码方面，大量中文文件显示为 mojibake，说明文件很可能曾以 UTF-8
保存、后被错误编码链处理。代码符号仍可读，但这是长期维护风险。

\begin{center}\rule{0.5\linewidth}{0.5pt}\end{center}

\section{第一部分：项目全景图}\label{ux7b2cux4e00ux90e8ux5206ux9879ux76eeux5168ux666fux56fe}

\section{ESP32 Draggy Racing Logger}\label{esp32-draggy-racing-logger}

\subsection{项目背景}\label{ux9879ux76eeux80ccux666f}

项目想用低成本器件做一个车载/摩托车数据终端：实时读取车辆运动状态，显示速度、姿态、G
力和性能测试结果，并让手机不安装 App 也能查看数据。

真正困难的不是``读出传感器''，而是车辆环境：

\begin{itemize}
\tightlist
\item
  发动机和路面引入高频振动。
\item
  MPU6050 有零偏、温漂和安装坐标误差。
\item
  GT-U7 在静止和弱信号时会速度漂移。
\item
  GPS 只有低频绝对速度，IMU 有高频响应但积分漂移。
\item
  Web、OLED、采样、录音和存储都跑在同一个 Arduino 主循环中。
\end{itemize}

\subsection{项目目标}\label{ux9879ux76eeux76eeux6807}

当前代码试图实现：

\begin{itemize}
\tightlist
\item
  200 Hz IMU 采样与滤波。
\item
  10 Hz GPS 解析、静止抑制和性能计时。
\item
  GPS-IMU 速度与 G 力的轻量融合。
\item
  摩托车倾角、弯道、翘头/翘尾和峰值统计。
\item
  OLED 本地显示。
\item
  ESP32 SoftAP + HTTP Web 仪表盘。
\item
  SPIFFS 低码率 WAV 录音。
\item
  SD 卡 10 Hz CSV 行驶日志。
\item
  与 PCB 配套的参数化 OpenSCAD 外壳。
\end{itemize}

必须区分：SD 记录代码存在，但当前启动流程没有初始化 SD
卡，所以默认不会真正记录。

\subsection{技术栈}\label{ux6280ux672fux6808}

{\def\LTcaptype{none} % do not increment counter
\begin{longtable}[]{@{}ll@{}}
\toprule\noalign{}
类别 & 使用内容 \\
\midrule\noalign{}
\endhead
\bottomrule\noalign{}
\endlastfoot
MCU & ESP32-WROOM-32 / PlatformIO \passthrough{\lstinline!esp32dev!} \\
框架 & Arduino Framework \\
语言 & C++、HTML/CSS/JavaScript、OpenSCAD \\
传感器 & MPU6050、GT-U7 GPS、模拟麦克风 \\
显示 & SSD1306 128x64 OLED \\
总线 & 双 I2C、UART2、HSPI、ADC1 \\
网络 & Wi-Fi SoftAP、HTTP \passthrough{\lstinline!WebServer!} \\
存储 & SPIFFS、SD/FAT CSV \\
数据格式 & NMEA、JSON、CSV、PCM WAV \\
构建 & PlatformIO \\
机械 & Gerber、OpenSCAD、STL \\
\end{longtable}
}

\subsection{项目一句话总结}\label{ux9879ux76eeux4e00ux53e5ux8bddux603bux7ed3}

基于 ESP32 构建低成本车载多传感器性能记录仪，融合 200 Hz IMU 与 10 Hz
GPS，实现姿态/G 力分析、性能计时、Web 实时可视化、录音与可扩展数据记录。

\subsection{项目目录结构分析}\label{ux9879ux76eeux76eeux5f55ux7ed3ux6784ux5206ux6790}

\begin{lstlisting}
esp32-imu/
├── src/                         固件主源码
│   ├── main.cpp                 系统装配、调度、OLED、HTTP API
│   ├── mpu6050.*                IMU 驱动、滤波、基础姿态
│   ├── gps_module.*             GPS、性能测试、速度融合、会话持久化
│   ├── performance_analyzer.*   G 力融合、车辆行为、摩托车分析
│   ├── audio_recorder.*         ADC 定时采样、WAV、自动录音
│   └── sd_logger.*              SD 初始化与 CSV 日志
├── data/index.html              SPIFFS Web 仪表盘
├── cad/                         PCB 外壳源码与 STL
├── platformio.ini               板卡、串口、依赖、SPIFFS 配置
├── README.md                    早期项目说明，已落后于源码
├── BUILD_GUIDE.md               PlatformIO 构建与排错说明
├── GPS_*.md                     GPS 接线、升级与优化思路
├── GT-U7_*.md                   GT-U7 优化总结，部分内容未与代码同步
├── INTERVIEW_NOTES.md           技术复盘素材，不是运行事实来源
├── Gerber_*.zip                 PCB 制造文件
├── wiring_diagram.txt           接线说明，接近当前双 I2C 方案
├── include/ lib/ test/          PlatformIO 模板目录，暂无自有实现/测试
└── interview/                   LaTeX/PDF 派生文档，不属于固件运行链
\end{lstlisting}

PCB 和外壳是同一产品化尝试的一部分，不应误识别为独立软件项目。Gerber
表明存在双层
PCB、丝印、阻焊、钻孔和飞针测试数据；无法仅从当前工作区确定实际焊接版本和最终
BOM。

\begin{center}\rule{0.5\linewidth}{0.5pt}\end{center}

\section{第二部分：系统架构设计}\label{ux7b2cux4e8cux90e8ux5206ux7cfbux7edfux67b6ux6784ux8bbeux8ba1}

\section{系统架构}\label{ux7cfbux7edfux67b6ux6784}

\subsection{总体架构}\label{ux603bux4f53ux67b6ux6784}

\begin{lstlisting}
                   +---------------- ESP32 ----------------+
MPU6050 -- I2C0 -->| 采样/均值/IIR --> 姿态与G力 --------+ |
                   |                                      | |
GT-U7 ---- UART2 ->| NMEA解析/质量判断/速度滤波 --------+ | |
                   |                                    | | |
Mic ------ ADC1 -->| Timer ISR --> WAV缓冲 --> SPIFFS   | | |
                   |                                    v v |
SD Card --- HSPI <-| CSV Logger <--- 融合/性能/摩托分析器  |
                   |                    |                 |
OLED ----- I2C1 <--| 本地显示 <---------+                 |
                   |                    v                 |
Phone <-- Wi-Fi ---| WebServer --> JSON API --> index.html|
                   +--------------------------------------+
\end{lstlisting}

\subsection{为什么这样拆分}\label{ux4e3aux4ec0ux4e48ux8fd9ux6837ux62c6ux5206}

\begin{itemize}
\tightlist
\item
  \passthrough{\lstinline!main.cpp!} 是 composition
  root：拥有硬件实例、启动顺序和时间调度。
\item
  驱动/算法按数据源拆分：IMU、GPS、音频、SD 各自维护状态。
\item
  \passthrough{\lstinline!performance\_analyzer!}
  位于传感器之上，读取全局 IMU/GPS 数据形成派生指标。
\item
  OLED 和 HTTP 是两个输出适配器，共享同一批全局状态。
\end{itemize}

这套拆法适合原型开发，新增功能快；代价是模块通过
\passthrough{\lstinline!extern!}
全局结构体耦合，调用顺序就是隐含的数据一致性协议。

\subsection{数据链路}\label{ux6570ux636eux94feux8def}

\subsubsection{运动数据链}\label{ux8fd0ux52a8ux6570ux636eux94fe}

\begin{lstlisting}
MPU6050 raw acceleration/gyro
        |
        v
10点移动平均 + 一阶IIR
        |
        +--> Roll/Pitch
        +--> G力/驾驶状态
        +--> 摩托倾角互补滤波
        +--> IMU速度积分
                         \
GT-U7 NMEA --> GPS速度 ---> 轻量互补修正 --> 性能指标/Web/OLED/CSV
\end{lstlisting}

GPS 提供长期不漂移的绝对速度；IMU
提供短期快速变化。当前实现不是严格卡尔曼滤波，而是偏置校准、死区、积分和
GPS 比例拉回。

\subsubsection{展示链}\label{ux5c55ux793aux94fe}

\begin{lstlisting}
全局状态结构体
   +--> OLED: 20 Hz，五种页面
   +--> /data: Web主仪表盘
   +--> /moto: 摩托数据
   +--> /audio: 录音状态
   +--> Serial: 5 Hz简报 / 2 s详细信息
\end{lstlisting}

\subsubsection{存储链}\label{ux5b58ux50a8ux94fe}

\begin{lstlisting}
Mic ADC --> 8-bit/8 kHz PCM --> 512 B buffer --> SPIFFS WAV
GPS+IMU+Fusion --> RideDataPoint --> CSV line --> SD card
Browser samples --> JavaScript memory --> client-side CSV/JSON export
\end{lstlisting}

这里实际上有三套``记录''：SPIFFS 音频、SD
原始/融合数据、浏览器内存记录。它们没有统一会话 ID 和时间基准。

\begin{center}\rule{0.5\linewidth}{0.5pt}\end{center}

\section{第三部分：项目运行流程}\label{ux7b2cux4e09ux90e8ux5206ux9879ux76eeux8fd0ux884cux6d41ux7a0b}

\section{上电初始化}\label{ux4e0aux7535ux521dux59cbux5316}

\begin{lstlisting}
Reset
  |
  v
Serial 115200 + BOOT输入
  |
  v
双I2C 400 kHz初始化 -> 扫描设备
  |
  v
MPU6050 -> GPS -> 性能分析 -> 摩托模式 -> 音频
  |
  v
SD初始化跳过（当前测试模式）
  |
  v
OLED -> SPIFFS -> Wi-Fi SoftAP -> HTTP路由
  |
  v
显示连接信息，进入 loop()
\end{lstlisting}

值得注意：\passthrough{\lstinline!initAudioRecorder()!} 在
\passthrough{\lstinline!SPIFFS.begin(true)!} 之前读取 SPIFFS
容量。首次启动或文件系统异常时，这个顺序不稳妥，应该先挂载文件系统。

\section{主循环调度}\label{ux4e3bux5faaux73afux8c03ux5ea6}

项目没有 FreeRTOS 显式任务，也没有统一 scheduler，而是
\passthrough{\lstinline!millis()!} 驱动的协作式时间片：

\begin{lstlisting}
loop()
  |
  +--> server.handleClient()
  +--> BOOT按键与串口命令
  |
  +--> 每5 ms:
  |      ReadMPU6050()
  |      fuseIMUSpeed()
  |
  +--> 每100 ms:
  |      updateGPSData()
  |      checkAutoLogging()
  |      optional logDataPoint()
  |
  +--> 每次循环:
  |      updateAudioRecorder()
  |
  +--> 每50 ms:
         updateFusedData()
         update OLED
         print serial
\end{lstlisting}

优点是易读、没有任务同步成本。问题是任何阻塞操作都会破坏所有``频率承诺''。HTTP
文件传输、GPS 重连 \passthrough{\lstinline!delay()!}、OLED I2C、SD flush
和串口输出都可能让 5 ms/125 us 级任务失约。

\section{中断流程}\label{ux4e2dux65adux6d41ux7a0b}

唯一明确的用户中断路径是音频定时器：

\begin{lstlisting}
125 us timer
   |
   v
ISR: adc1_get_raw()
   |
   +--> lastSample = sample
   +--> sampleReady = true
             |
             v
主循环 processSample()
   |
   +--> 12 bit转8 bit
   +--> 更新音量
   +--> 写入512 B缓冲
   +--> 缓冲满后写SPIFFS
\end{lstlisting}

这不是可靠的生产者-消费者队列。若主循环在两个中断之间没有消费，\passthrough{\lstinline!lastSample!}
被覆盖，只保留最新样本。

\section{通信流程}\label{ux901aux4fe1ux6d41ux7a0b}

\subsection{GPS}\label{gps}

\begin{lstlisting}
UART2 bytes
  -> TinyGPS++ encode
  -> location/speed/time/satellite/HDOP
  -> 合理范围与跳变检查
  -> 速度20点窗口、方差、自适应滤波、死区
  -> GPS_DATA
\end{lstlisting}

超过 3 秒无可解析数据会标记丢失；每 10
秒发送热启动和恢复配置。恢复配置将更新周期设置为 5 Hz，与启动时 10 Hz
不一致。

\subsection{Web}\label{web}

ESP32 创建 \passthrough{\lstinline!ESP32-IMU!} 热点，密码硬编码为
\passthrough{\lstinline!12345678!}，地址通常为
\passthrough{\lstinline!192.168.4.1!}。

{\def\LTcaptype{none} % do not increment counter
\begin{longtable}[]{@{}ll@{}}
\toprule\noalign{}
路由 & 作用 \\
\midrule\noalign{}
\endhead
\bottomrule\noalign{}
\endlastfoot
\passthrough{\lstinline!/!} & 从 SPIFFS 流式返回
\passthrough{\lstinline!index.html!} \\
\passthrough{\lstinline!/data!} & 通用速度、GPS、G 力、性能 JSON \\
\passthrough{\lstinline!/moto!} & 摩托倾角、事件和峰值 JSON \\
\passthrough{\lstinline!/audio!} & 录音状态 JSON \\
\passthrough{\lstinline!/audio/control?action=...!} &
开始、停止、暂停、恢复、删除、自动模式、列表 \\
\passthrough{\lstinline!/audio/download?file=...!} & 下载 WAV \\
\end{longtable}
}

网页每 150 ms 请求 \passthrough{\lstinline!/data!}，每 100 ms 请求
\passthrough{\lstinline!/moto!}。这是轮询架构，简单但请求频繁，且
\passthrough{\lstinline!/moto!} 请求本身会再次调用
\passthrough{\lstinline!updateMotorcycleData()!}，使算法更新时间受浏览器请求节奏影响。

\section{状态机}\label{ux72b6ux6001ux673a}

\subsection{音频录制}\label{ux97f3ux9891ux5f55ux5236}

\begin{lstlisting}
IDLE/STOPPED -- speed >= 5 --> ACTIVE
ACTIVE -- speed < 2 for 5s --> STOPPED
ACTIVE -- pause ------------> PAUSED
PAUSED -- resume -----------> ACTIVE
ACTIVE -- 300s -------------> STOPPED
\end{lstlisting}

\passthrough{\lstinline!RECORDING\_WAITING!}
被定义但当前逻辑基本未使用。

\subsection{性能测试}\label{ux6027ux80fdux6d4bux8bd5}

\begin{itemize}
\tightlist
\item
  加速：速度超过 5 km/h 自动开始，穿越 60/100/200 km/h 时记录时间。
\item
  制动：速度达到约 95 km/h 后进入条件，降到 5 km/h 结束。
\item
  圈速字段和函数声明存在，但完整实现无法从当前工作区确定。
\end{itemize}

\begin{center}\rule{0.5\linewidth}{0.5pt}\end{center}

\section{第四部分：核心模块解析}\label{ux7b2cux56dbux90e8ux5206ux6838ux5fc3ux6a21ux5757ux89e3ux6790}

\section{\texorpdfstring{\texttt{main.cpp}：系统装配与输出层}{main.cpp：系统装配与输出层}}\label{maincppux7cfbux7edfux88c5ux914dux4e0eux8f93ux51faux5c42}

职责：

\begin{itemize}
\tightlist
\item
  创建两条 \passthrough{\lstinline!TwoWire!} 总线、OLED、WebServer。
\item
  固定初始化顺序和主循环频率。
\item
  把全局数据编码为 OLED 页面和 JSON。
\item
  管理 BOOT 按键及串口模式。
\end{itemize}

关键设计是多速率协作循环。最大问题是文件已超过千行，HTTP、OLED、调度和业务拼在一起，修改一个输出字段容易影响实时路径。

\section{\texorpdfstring{\texttt{mpu6050}：高振动车辆 IMU
前端}{mpu6050：高振动车辆 IMU 前端}}\label{mpu6050ux9ad8ux632fux52a8ux8f66ux8f86-imu-ux524dux7aef}

核心结构 \passthrough{\lstinline!MPU6050\_DATA!}
同时保存原始和滤波后的加速度、角速度、Roll/Pitch。

处理链：

\begin{lstlisting}
Adafruit getEvent()
 -> 10样本滑动平均
 -> alpha=0.1 一阶IIR
 -> 加速度倾角
\end{lstlisting}

量程为 ±8 g、±500 deg/s，DLPF 为 94 Hz。选择 94 Hz
是为了保留快速车辆动作，再依赖软件滤波压制振动。

注意：

\begin{itemize}
\tightlist
\item
  \passthrough{\lstinline!CalculateAttitude()!} 的基础 Roll/Pitch
  仅使用加速度计。
\item
  \passthrough{\lstinline!Yaw!} 固定为 0。
\item
  陀螺数据在 Adafruit API 中是 rad/s，但摩托互补滤波直接按 deg/s
  积分的迹象很强；若未转换，倾角动态项会缩小约 57.3
  倍。需要实机与库单位确认后修正。
\item
  初始滤波状态为 0，会有启动爬升过程。
\end{itemize}

\section{\texorpdfstring{\texttt{gps\_module}：GPS、速度与性能计时}{gps\_module：GPS、速度与性能计时}}\label{gps_modulegpsux901fux5ea6ux4e0eux6027ux80fdux8ba1ux65f6}

核心结构：

\begin{itemize}
\tightlist
\item
  \passthrough{\lstinline!GPS\_DATA!}：位置、时间、速度、信号质量、融合状态。
\item
  \passthrough{\lstinline!PERFORMANCE\_DATA!}：加速、制动、圈速和会话数据。
\end{itemize}

主要策略：

\begin{itemize}
\tightlist
\item
  仅输出 RMC/GGA，降低 9600 bps 下 NMEA 堵塞风险。
\item
  开启 SBAS/WAAS。
\item
  对位置做经纬度范围和 1 km 单帧跳变检查。
\item
  对速度做 20 点窗口、标准差、卫星/HDOP 动态阈值、分段 IIR 和最终死区。
\item
  静止时收集 200 个样本估计 IMU 前向偏置。
\item
  IMU 积分速度每次 GPS 新帧向 GPS 拉回 10\%。
\end{itemize}

设计意图正确，但实现中 \passthrough{\lstinline!gps\_data.speed\_ms!}
很少被明确同步为当前 GPS 速度，而校准与静止归零依赖它；应统一只使用
\passthrough{\lstinline!speed\_mps!} 或明确维护该字段。

\section{\texorpdfstring{\texttt{performance\_analyzer}：派生指标与摩托语义}{performance\_analyzer：派生指标与摩托语义}}\label{performance_analyzerux6d3eux751fux6307ux6807ux4e0eux6469ux6258ux8bedux4e49}

\passthrough{\lstinline!FUSED\_PERFORMANCE\_DATA!} 保存当前速度/G
力、峰值、行为状态和质量分数。

主要功能：

\begin{itemize}
\tightlist
\item
  组合 IMU 加速度与 GPS 速度差分。
\item
  检测加速、制动、转弯。
\item
  用车辆质量、滚阻、风阻估算瞬时功率。
\item
  用互补滤波估计摩托倾角和俯仰。
\item
  检测弯道、翘头、翘尾，估算弯道半径。
\end{itemize}

这里把``信号处理''和``车辆业务语义''放在同一模块，原型阶段方便，但未来应拆成：

\begin{lstlisting}
Sensor Fusion
Vehicle Frame Transform
Event Detector
Session Statistics
\end{lstlisting}

硬编码偏置 \passthrough{\lstinline!-0.47/+0.5/-0.48!} 和假定坐标
\passthrough{\lstinline!X=垂直、Y=侧向、Z=前向!}
是整个分析可信度的基础。换安装方向或换板后必须重新标定。

\section{\texorpdfstring{\texttt{audio\_recorder}：低码率事件录音}{audio\_recorder：低码率事件录音}}\label{audio_recorderux4f4eux7801ux7387ux4e8bux4ef6ux5f55ux97f3}

配置为单声道、8 kHz、8 bit PCM，每秒约 8 KB，5 分钟约 2.4
MB，不含文件系统开销。

WAV
头在开始时写占位值，停止时回写文件大小。这种方式简单，但异常断电会留下头部未完成的文件。

最大架构问题是 ISR 只有一个样本槽，不是
DMA/环形缓冲。录音数据完整性无法保证。

\section{\texorpdfstring{\texttt{sd\_logger}：行驶
CSV}{sd\_logger：行驶 CSV}}\label{sd_loggerux884cux9a76-csv}

SD 使用 HSPI：

\begin{lstlisting}
MISO GPIO25
MOSI GPIO13
SCK  GPIO14
CS   GPIO15
\end{lstlisting}

选择 GPIO25 替代 GPIO12，是为了规避 ESP32 启动绑带引脚风险。日志以 10 Hz
写 CSV，每 100 点 flush 一次，在数据安全和写入开销之间折中。

当前 \passthrough{\lstinline!setup()!} 没有调用
\passthrough{\lstinline!initSDCard()!}，因此
\passthrough{\lstinline!checkAutoLogging()!} 会因
\passthrough{\lstinline!initialized=false!} 直接返回。

\section{\texorpdfstring{\texttt{data/index.html}：无 App Web
仪表盘}{data/index.html：无 App Web 仪表盘}}\label{dataindexhtmlux65e0-app-web-ux4eeaux8868ux76d8}

单文件前端包含：

\begin{itemize}
\tightlist
\item
  实时速度/G 力曲线。
\item
  性能、GPS、摩托状态。
\item
  浏览器端会话记录与 CSV/JSON 导出。
\item
  录音控制逻辑。
\end{itemize}

优点是部署简单。缺点是 HTML 已很大、编码损坏明显、无构建/校验、API
字段靠手写约定。浏览器记录只存在内存，刷新即丢失。

\section{\texorpdfstring{\texttt{cad/pcb1\_enclosure.scad}：参数化产品外壳}{cad/pcb1\_enclosure.scad：参数化产品外壳}}\label{cadpcb1_enclosurescadux53c2ux6570ux5316ux4ea7ux54c1ux5916ux58f3}

外壳从 Gerber 轮廓推导出 78.867 x 70.866 mm、圆角 5.08 mm 的 PCB：

\begin{itemize}
\tightlist
\item
  底壳用周边支撑环和卡扣固定无安装孔 PCB。
\item
  上盖使用局部器件避让腔。
\item
  四角外置螺钉耳避免占用板内空间。
\item
  提供通风槽、边缘开口、装配/爆炸/导出模式。
\end{itemize}

器件高度和接口位置是从丝印/飞针数据推断的，不是装配实测。打印前必须实测校核。

\begin{center}\rule{0.5\linewidth}{0.5pt}\end{center}

\section{第五部分：技术知识提取}\label{ux7b2cux4e94ux90e8ux5206ux6280ux672fux77e5ux8bc6ux63d0ux53d6}

\section{双 I2C}\label{ux53cc-i2c}

\subsection{为什么需要它}\label{ux4e3aux4ec0ux4e48ux9700ux8981ux5b83}

MPU6050 与 OLED 本可共享地址不同的总线，但 OLED 整帧刷新会占用 I2C
时间。分成两条控制器可以隔离显示流量与传感器采样，也方便独立排错。

\subsection{项目中的实际应用}\label{ux9879ux76eeux4e2dux7684ux5b9eux9645ux5e94ux7528}

\begin{itemize}
\tightlist
\item
  I2C0：GPIO21/22，MPU6050。
\item
  I2C1：GPIO18/19，SSD1306。
\item
  两者均配置 400 kHz。
\end{itemize}

\subsection{常见 Bug 与调试}\label{ux5e38ux89c1-bug-ux4e0eux8c03ux8bd5}

\begin{itemize}
\tightlist
\item
  文档接线漂移是当前首要风险。
\item
  地址应扫描确认：MPU 0x68/0x69，OLED 0x3C/0x3D。
\item
  线长、上拉、电源噪声会导致间歇 NACK。
\item
  用现有 \passthrough{\lstinline!scanI2CDevices()!}
  先区分接线问题与驱动问题。
\end{itemize}

替代方案是共享一条 I2C，节省引脚；对本项目而言隔离总线更适合实时采样。

\section{UART 与 NMEA}\label{uart-ux4e0e-nmea}

GT-U7 通过 UART2 输入 NMEA。项目只保留 RMC 与 GGA，因为 10 Hz、9600 bps
时全量语句容易占满串口带宽。

常见坑：

\begin{itemize}
\tightlist
\item
  模块不一定兼容所有 PMTK 命令。
\item
  发送 10 Hz 配置不代表模块确实进入 10 Hz。
\item
  掉线恢复又设置 5 Hz，导致运行状态不一致。
\item
  应统计每秒完整 RMC/GGA 帧数，而不是只相信配置命令。
\end{itemize}

\section{采样与数字滤波}\label{ux91c7ux6837ux4e0eux6570ux5b57ux6ee4ux6ce2}

项目使用了三类低成本滤波：

\begin{itemize}
\tightlist
\item
  硬件 DLPF：在传感器内部先限制带宽。
\item
  滑动平均：压制随机噪声，但引入窗口延迟。
\item
  一阶 IIR：O(1) 状态，适合 MCU，但有相位延迟。
\end{itemize}

常见误区是把``200 Hz
调用周期''当成严格采样率。当前系统是协作循环，实际采样间隔会抖动，且代码给融合器的
\passthrough{\lstinline!dt!} 固定为 5 ms，没有用真实时间差。

调试时应同时记录：

\begin{itemize}
\tightlist
\item
  实际 \passthrough{\lstinline!micros()!} 间隔分布。
\item
  原始值、滑动平均值、IIR 值。
\item
  静止 RMS、阶跃响应和相位延迟。
\end{itemize}

\section{GPS 静止漂移}\label{gps-ux9759ux6b62ux6f02ux79fb}

GPS
低速时噪声可能大于真实速度。项目不是简单小于阈值清零，而是综合窗口方差、连续稳定次数、卫星数和
HDOP。

这能改善显示，但阈值滤波也会吃掉真实低速起步，使 0-100
起点变晚。性能计时最好保留原始 Doppler
速度和滤波速度两路，分别服务计时与 UI。

\section{GPS-IMU 融合}\label{gps-imu-ux878dux5408}

当前思路：

\begin{lstlisting}
IMU: 高频预测，但漂移
GPS: 低频校正，但延迟和低速噪声
\end{lstlisting}

项目采用简化互补校正而非
EKF，优点是参数直观、计算低。缺点是没有显式状态协方差、时间戳对齐、姿态重力补偿和传感器延迟模型。

真正升级时，优先实现一维状态
\passthrough{\lstinline![速度, 加速度偏置]!} 的 Kalman
filter，而不是直接上完整 9/15 状态导航 EKF。

\section{ADC、Timer 与音频}\label{adctimer-ux4e0eux97f3ux9891}

定时器中断可以建立稳定采样节拍，但 ISR
应只搬运数据，不能依赖单槽共享变量承载连续流。

更合适的方案：

\begin{itemize}
\tightlist
\item
  ESP32 I2S ADC/DMA。
\item
  ISR 写入 lock-free ring buffer。
\item
  双缓冲，由任务批量写文件。
\end{itemize}

此外 8-bit unsigned PCM 的直流中心假设是 128；模拟麦克风偏置和 ADC
非线性应实测。

\section{SPI、SD 与文件系统}\label{spisd-ux4e0eux6587ux4ef6ux7cfbux7edf}

SD 写入延迟具有长尾，\passthrough{\lstinline!flush()!}
可能阻塞几十毫秒。若和 200 Hz 采样同循环，必须缓冲解耦。

SPIFFS 适合少量静态资源和小文件，不适合高频持续音频写入。LittleFS 或 SD
更适合长期录音；异常断电还需要临时文件和恢复策略。

\section{Wi-Fi 与 HTTP}\label{wi-fi-ux4e0e-http}

SoftAP 让手机直接连接，现场无需路由器，这是很实用的产品决策。

当前轮询方案的常见问题：

\begin{itemize}
\tightlist
\item
  多客户端会放大 CPU 和 JSON 序列化开销。
\item
  同步 \passthrough{\lstinline!WebServer!} 的流式下载可能阻塞采样。
\item
  \passthrough{\lstinline!sprintf!} 固定缓冲区依赖字段长度，需用
  \passthrough{\lstinline!snprintf!}。
\item
  密码硬编码，适合原型但不适合发布。
\end{itemize}

替代方案：Server-Sent Events 或 WebSocket
推送，并将静态资源缓存与实时数据通道分离。

\begin{center}\rule{0.5\linewidth}{0.5pt}\end{center}

\section{第六部分：设计决策分析}\label{ux7b2cux516dux90e8ux5206ux8bbeux8ba1ux51b3ux7b56ux5206ux6790}

\section{双 I2C
而不是共享总线}\label{ux53cc-i2c-ux800cux4e0dux662fux5171ux4eabux603bux7ebf}

背景：显示更新与高频 IMU 读取争用总线。

优点：故障隔离、时序可控、扫描清晰。缺点：多占两根 GPIO，接线复杂。当前
ESP32 引脚尚可，选择合理。

\section{协作式调度而不是 FreeRTOS
任务}\label{ux534fux4f5cux5f0fux8c03ux5ea6ux800cux4e0dux662f-freertos-ux4efbux52a1}

背景：Arduino 原型需要快速迭代。

优点：控制流直观、调试容易。缺点：所有模块共享阻塞预算，无法保证 200 Hz
IMU 和 8 kHz 音频。

当前方案适合功能验证；当音频、SD 和网络同时开启时，应迁移为任务和队列：

\begin{lstlisting}
Core/Task Sensor -> Queue -> Fusion
Core/Task IO     -> SD/SPIFFS/Web
\end{lstlisting}

\section{多级滤波而不是单一滤波器}\label{ux591aux7ea7ux6ee4ux6ce2ux800cux4e0dux662fux5355ux4e00ux6ee4ux6ce2ux5668}

车辆振动覆盖宽频段，单一强低通会造成明显延迟。硬件 DLPF、移动平均和 IIR
分层容易调参，但当前级联实际带宽没有通过测试数据验证。

替代方案包括双二阶 Butterworth、陷波器、Savitzky-Golay
或频域分析后定制滤波。应先采集原始数据做 PSD，再决定，而不是继续堆滤波。

\section{互补滤波而不是
EKF}\label{ux4e92ux8865ux6ee4ux6ce2ux800cux4e0dux662f-ekf}

当前传感器精度有限，且原型更需要可解释性。互补滤波成本低、故障容易定位。

缺点是重力投影、安装误差和时间同步需要手工处理。等硬件、坐标系和时间戳稳定后，再引入一维
KF 或姿态 EKF。

\section{阈值状态机}\label{ux9608ux503cux72b6ux6001ux673a}

自动录音、SD 日志、加速/制动、弯道检测都采用阈值 + 延时/滞回。

优点是行为可解释；缺点是阈值分散在多个文件，且录音与 SD
的起停参数不一致：

\begin{itemize}
\tightlist
\item
  录音停止延时 5 秒。
\item
  SD 停止延时 30 秒。
\end{itemize}

应集中成
\passthrough{\lstinline!VehicleSessionConfig!}，并用一次会话状态驱动所有记录器。

\section{浏览器端导出}\label{ux6d4fux89c8ux5668ux7aefux5bfcux51fa}

浏览器记录减轻 ESP32
存储压力，开发方便。但手机锁屏、刷新、断连都会丢数据，不能替代设备端可靠日志。它更适合作为``即时观察与临时导出''。

\begin{center}\rule{0.5\linewidth}{0.5pt}\end{center}

\section{第七部分：调试与踩坑记录}\label{ux7b2cux4e03ux90e8ux5206ux8c03ux8bd5ux4e0eux8e29ux5751ux8bb0ux5f55}

\section{调试记录
1：文档与固件接线冲突}\label{ux8c03ux8bd5ux8bb0ux5f55-1ux6587ux6863ux4e0eux56faux4ef6ux63a5ux7ebfux51b2ux7a81}

问题：按旧 README 接线可能找不到设备。

原因：项目从共享 I2C 演进到双 I2C，文档未统一。

定位：以 \passthrough{\lstinline!main.cpp!} 宏和启动扫描输出为准。

解决：维护唯一
\passthrough{\lstinline!HARDWARE\_PINOUT.md!}，其他文档只链接，不复制引脚表。

经验：硬件配置必须单一事实源。

\section{调试记录
2：音频丢采样}\label{ux8c03ux8bd5ux8bb0ux5f55-2ux97f3ux9891ux4e22ux91c7ux6837}

问题：WAV 时长、音质或样本数异常。

原因：8 kHz ISR 覆盖单个
\passthrough{\lstinline!lastSample!}；主循环无法保证每 125 us 消费。

定位：比较墙钟时长与
\passthrough{\lstinline!samples\_recorded / 8000!}，统计覆盖次数。

解决：I2S DMA 或环形缓冲，文件写入放到独立任务。

\section{调试记录 3：SD
功能看似存在但不工作}\label{ux8c03ux8bd5ux8bb0ux5f55-3sd-ux529fux80fdux770bux4f3cux5b58ux5728ux4f46ux4e0dux5de5ux4f5c}

问题：速度超过阈值仍不生成 CSV。

原因：\passthrough{\lstinline!initSDCard()!} 在
\passthrough{\lstinline!setup()!}
被注释，\passthrough{\lstinline!sd\_status.initialized!} 始终为 false。

定位：查看启动日志是否打印 SD 容量。

解决：恢复初始化前先解决原注释所说的 I2C/引脚冲突，并增加 Web/串口 SD
状态。

\section{调试记录
4：倾角动态响应可能错误}\label{ux8c03ux8bd5ux8bb0ux5f55-4ux503eux89d2ux52a8ux6001ux54cdux5e94ux53efux80fdux9519ux8bef}

问题：陀螺贡献过弱，倾角主要跟随加速度计。

原因：Adafruit 陀螺输出通常为 rad/s，而摩托代码按 deg/s 注释直接积分。

定位：固定角速度转动，比较原始 gyro、积分角与实际角度。

解决：在驱动边界统一单位，字段命名加入
\passthrough{\lstinline!\_rad\_s!} 或
\passthrough{\lstinline!\_deg\_s!}。

\section{调试记录 5：GPS
模式不一致}\label{ux8c03ux8bd5ux8bb0ux5f55-5gps-ux6a21ux5f0fux4e0dux4e00ux81f4}

问题：掉线恢复后刷新率下降，性能测试响应变化。

原因：启动发送 100 ms 更新命令，恢复发送 200 ms 更新命令。

解决：抽取
\passthrough{\lstinline!configureGPS()!}，启动和恢复调用同一配置；读取实际帧率确认结果。

\section{调试记录 6：启动顺序与
SPIFFS}\label{ux8c03ux8bd5ux8bb0ux5f55-6ux542fux52a8ux987aux5e8fux4e0e-spiffs}

问题：音频模块初始化时容量可能为 0 或状态不可靠。

原因：音频初始化早于 \passthrough{\lstinline!SPIFFS.begin(true)!}。

解决：先初始化存储，再初始化依赖存储的音频和 Web。

\section{调试记录 7：融合字段与 API
不一致}\label{ux8c03ux8bd5ux8bb0ux5f55-7ux878dux5408ux5b57ux6bb5ux4e0e-api-ux4e0dux4e00ux81f4}

问题：内部已经运行速度融合，但网页 \passthrough{\lstinline!imuSpeed!}
永远为 0，\passthrough{\lstinline!speed!} 仍来自 GPS。

原因：算法、状态结构和 JSON 输出没有同时更新。

解决：定义明确字段：\passthrough{\lstinline!gps\_speed!}、\passthrough{\lstinline!imu\_predicted\_speed!}、\passthrough{\lstinline!fused\_speed!}，并建立
API schema 测试。

\section{调试记录 8：HTTP
文件下载路径}\label{ux8c03ux8bd5ux8bb0ux5f55-8http-ux6587ux4ef6ux4e0bux8f7dux8defux5f84}

问题：下载接口可能访问任意 SPIFFS 已知路径，而不仅是录音。

原因：只补 \passthrough{\lstinline!/!} 和检查文件存在，没有限定
\passthrough{\lstinline!/ride\_*.wav!}。

解决：白名单文件名格式，拒绝 \passthrough{\lstinline!..!}、非 WAV
和非录音前缀。

\section{调试记录
9：字符串与缓冲区}\label{ux8c03ux8bd5ux8bb0ux5f55-9ux5b57ux7b26ux4e32ux4e0eux7f13ux51b2ux533a}

问题：JSON/文本字段增长后可能溢出固定缓冲区。

原因：使用 \passthrough{\lstinline!sprintf!} 而非长度受限接口。

解决：改为 \passthrough{\lstinline!snprintf!} 并检查返回值；长期使用
ArduinoJson 流式序列化。

\section{调试记录
10：编码损坏}\label{ux8c03ux8bd5ux8bb0ux5f55-10ux7f16ux7801ux635fux574f}

问题：中文日志、网页标签和文档显示乱码，甚至可能破坏 HTML 标签或 C++
字符串。

原因：文件编码链不一致。

解决：全仓统一 UTF-8，无 BOM；添加
\passthrough{\lstinline!.editorconfig!}，在 CI 中做 UTF-8 解码和
HTML/C++ 基础检查。

\begin{center}\rule{0.5\linewidth}{0.5pt}\end{center}

\section{第八部分：工程实践总结}\label{ux7b2cux516bux90e8ux5206ux5de5ux7a0bux5b9eux8df5ux603bux7ed3}

\subsection{做得好的地方}\label{ux505aux5f97ux597dux7684ux5730ux65b9}

\begin{itemize}
\tightlist
\item
  模块按传感器和业务域拆分，原型意图清楚。
\item
  使用 \passthrough{\lstinline!millis()!} 而不是在主路径大量固定延时。
\item
  对 GPS 弱信号、静止漂移和跳点有工程化防护。
\item
  同时保留原始值、滤波值和峰值，便于调试。
\item
  通过 SoftAP + 单页 Web 降低现场使用门槛。
\item
  SD 选择 HSPI 和非 GPIO12 MISO，体现了对 ESP32 启动绑带的考虑。
\item
  OpenSCAD 参数化外壳把软件项目向可装配设备推进了一步。
\end{itemize}

\subsection{可以改进的地方}\label{ux53efux4ee5ux6539ux8fdbux7684ux5730ux65b9}

\subsubsection{代码组织}\label{ux4ee3ux7801ux7ec4ux7ec7}

\passthrough{\lstinline!main.cpp!} 应拆成：

\begin{lstlisting}
app_scheduler
web_api
oled_view
hardware_config
\end{lstlisting}

全局 \passthrough{\lstinline!extern!}
状态应逐步替换为只读快照或消息结构，至少明确谁写、谁读、更新频率和单位。

\subsubsection{配置管理}\label{ux914dux7f6eux7ba1ux7406}

引脚、采样率、阈值、车辆参数和 Wi-Fi 凭据分散在各文件。建议集中：

\begin{lstlisting}
board_config.h
sensor_config.h
vehicle_config.h
\end{lstlisting}

并把所有物理量写入单位后缀。

\subsubsection{日志}\label{ux65e5ux5fd7}

串口日志丰富，但缺乏级别和节流统一。应使用
\passthrough{\lstinline!LOGE/W/I/D!}，避免高频
\passthrough{\lstinline!Serial.printf!} 破坏实时性。

\subsubsection{错误处理}\label{ux9519ux8befux5904ux7406}

很多初始化失败后继续运行，这是原型的容错思路，但后续函数仍可能假定硬件有效。应维护模块健康状态，并在
UI 显示 degraded mode。

\subsubsection{测试策略}\label{ux6d4bux8bd5ux7b56ux7565}

当前 \passthrough{\lstinline!test/!}
只有模板说明，没有项目测试。最值得先补的是纯算法主机测试：

\begin{itemize}
\tightlist
\item
  Haversine 距离。
\item
  GPS 窗口与静止判定。
\item
  阈值穿越计时。
\item
  WAV 头尺寸。
\item
  CSV 格式。
\item
  互补滤波单位与步进响应。
\item
  状态机起停滞回。
\end{itemize}

硬件在环测试应回放已保存的 NMEA 与 IMU 数据，而不是每次都上车测试。

\subsubsection{构建与 CI}\label{ux6784ux5efaux4e0e-ci}

当前配置锁定库版本是好事。建议 CI 至少执行：

\begin{lstlisting}
pio run
pio test -e native
HTML/JSON schema sanity check
UTF-8 validation
\end{lstlisting}

本次分析环境无法执行构建，因为终端没有可用的
\passthrough{\lstinline!pio!}
命令；因此当前源码是否能完整编译无法从本次静态分析确定。

\begin{center}\rule{0.5\linewidth}{0.5pt}\end{center}

\section{第九部分：项目复盘}\label{ux7b2cux4e5dux90e8ux5206ux9879ux76eeux590dux76d8}

\subsection{最大的亮点}\label{ux6700ux5927ux7684ux4eaeux70b9}

不是简单堆传感器，而是围绕``低成本硬件在真实车辆环境中不可靠''做了多层处理：信号质量、静止抑制、漂移校准、融合、状态机和多终端显示。

\subsection{最大的难点}\label{ux6700ux5927ux7684ux96beux70b9}

多时间尺度和多误差源同时存在：

\begin{itemize}
\tightlist
\item
  音频 8 kHz。
\item
  IMU 200 Hz。
\item
  显示/融合 20 Hz。
\item
  GPS/日志 10 Hz。
\item
  Web 约 6.7/10 Hz。
\end{itemize}

而它们目前共享一个可能阻塞的主循环。真正的系统瓶颈是调度与数据一致性，不再只是算法。

\subsection{学到的知识}\label{ux5b66ux5230ux7684ux77e5ux8bc6}

\begin{itemize}
\tightlist
\item
  传感器数据``能读''与``能用于决策''之间隔着校准、坐标系、滤波、质量评估和时间同步。
\item
  低成本 GPS
  的极限不能只靠软件宣传突破，升级硬件有时比继续堆阈值更合理。
\item
  原型文档很容易落后于代码，尤其是引脚和单位。
\item
  记录系统必须考虑掉电、断连和阻塞长尾。
\item
  机械结构、接线和软件配置应被视为同一系统架构。
\end{itemize}

\subsection{如果重新设计}\label{ux5982ux679cux91cdux65b0ux8bbeux8ba1}

\begin{enumerate}
\def\labelenumi{\arabic{enumi}.}
\tightlist
\item
  先定义统一硬件配置、坐标系和单位。
\item
  使用 I2S DMA 采音，传感器和存储通过队列解耦。
\item
  建立单一 \passthrough{\lstinline!VehicleSession!}
  状态机，统一音频、CSV 和峰值会话。
\item
  所有数据携带 \passthrough{\lstinline!micros()!} 时间戳，融合使用真实
  \passthrough{\lstinline!dt!}。
\item
  保留 raw/filtered/fused 三层数据并定义稳定 API schema。
\item
  先做数据回放工具和 native 单元测试，再调算法。
\item
  用实测频谱和标定数据选择滤波器。
\item
  把网页通信改成 SSE/WebSocket。
\end{enumerate}

\subsection{下一步扩展}\label{ux4e0bux4e00ux6b65ux6269ux5c55}

优先级建议：

\begin{enumerate}
\def\labelenumi{\arabic{enumi}.}
\tightlist
\item
  修复编码、引脚文档和单位问题。
\item
  恢复可验证的 SD 日志，获得真实数据集。
\item
  修复音频采样架构。
\item
  增加时间戳和离线回放测试。
\item
  一维速度/偏置 Kalman filter。
\item
  升级 NEO-M9N 或更高频 GNSS；是否上 RTK 取决于实际计时精度目标。
\item
  增加 OTA、配置页面和会话下载。
\end{enumerate}

无法从当前工作区确定：现有 PCB 是否已经覆盖麦克风与 SD
模块、外壳是否完成实物试装、车辆参数是否经过标定、GPS
优化报告中的量化提升是否来自可复现实测。

\begin{center}\rule{0.5\linewidth}{0.5pt}\end{center}

\section{第十部分：个人知识地图}\label{ux7b2cux5341ux90e8ux5206ux4e2aux4ebaux77e5ux8bc6ux5730ux56fe}

\section{Embedded Knowledge Map}\label{embedded-knowledge-map}

\begin{lstlisting}
嵌入式系统
├── ESP32
│   ├── Arduino Framework
│   ├── GPIO与启动绑带
│   ├── 双I2C控制器
│   ├── UART2
│   ├── HSPI
│   ├── ADC1与硬件Timer
│   ├── SPIFFS
│   └── Wi-Fi SoftAP / HTTP
│
├── 传感器与信号处理
│   ├── MPU6050量程与DLPF
│   ├── 滑动平均
│   ├── 一阶IIR
│   ├── 互补滤波
│   ├── 零偏与安装校准
│   ├── 坐标系映射
│   ├── 重力补偿
│   └── 采样率、抖动与真实dt
│
├── GNSS
│   ├── NMEA RMC/GGA
│   ├── Doppler速度
│   ├── 卫星数与HDOP
│   ├── 静止漂移抑制
│   ├── 跳点过滤
│   ├── SBAS
│   ├── 多频率数据融合
│   └── GT-U7到高性能GNSS升级
│
├── 车辆与机器人数据
│   ├── 纵向/横向/垂直G
│   ├── 加速与制动计时
│   ├── 倾角与俯仰
│   ├── 弯道检测与半径估计
│   ├── 翘头/翘尾事件
│   ├── 峰值与会话统计
│   └── 简化功率估算
│
├── 数据记录
│   ├── WAV/PCM
│   ├── CSV
│   ├── JSON API
│   ├── 环形缓冲与DMA
│   ├── SD写入长尾
│   ├── 掉电一致性
│   └── 浏览器端临时记录
│
├── 软件架构
│   ├── 协作式调度
│   ├── 多速率任务
│   ├── 状态机与滞回
│   ├── 全局状态耦合
│   ├── 消息队列/生产者消费者
│   ├── 模块健康状态
│   └── API schema
│
├── 工程化
│   ├── PlatformIO依赖锁定
│   ├── Native算法测试
│   ├── 硬件在环与数据回放
│   ├── 日志分级
│   ├── UTF-8与文档单一事实源
│   └── CI构建与静态检查
│
└── 产品化
    ├── Gerber与PCB制造资料
    ├── 飞针测试数据
    ├── OpenSCAD参数化外壳
    ├── 装配公差与器件避让
    ├── Web免安装交互
    └── 硬件成本与精度边界
\end{lstlisting}

\subsection{长期维护入口}\label{ux957fux671fux7ef4ux62a4ux5165ux53e3}

半年后重新开始时，建议按这个顺序：

\begin{enumerate}
\def\labelenumi{\arabic{enumi}.}
\tightlist
\item
  先读本文的``阅读前事实''和``调试记录''。
\item
  核对实际硬件接线与 \passthrough{\lstinline!main.cpp!} 引脚。
\item
  执行 PlatformIO 构建并保存构建环境版本。
\item
  串口确认 I2C 扫描、GPS 帧率、SPIFFS 和 SD 状态。
\item
  先记录一份 raw 数据，再修改滤波或融合参数。
\item
  所有新的实测结论更新到本文，同时删除或标记过时文档中的冲突说法。
\end{enumerate}

\end{document}
